\documentclass{ctexart}
\newcommand{\keywords}[1]{\par\noindent\textbf{关键词：}#1}
\begin{document}
\begin{abstract}
本文针对有限资源配置问题建立线性规划基线，并在相同约束下比较两种方案。模型给出可复算的目标值、约束余量和误差检验。
\end{abstract}
\keywords{资源配置；线性规划；方案比较}
\section{问题重述}
给定三类资源及其容量，在预算约束下确定两项决策变量，使总收益最大。
\section{问题分析}
由于目标和约束均可写成线性形式，先建立线性规划基线；再扰动容量，观察最优方案是否切换。
\section{模型假设}
假设单位收益在研究区间内不变，因此收益可由决策量线性叠加。
\section{符号说明}
$x_i$ 表示第 $i$ 项决策量，$c_i$ 表示单位收益，$b_j$ 表示第 $j$ 类资源上限。
\section{分问题模型建立与求解}
目标函数为 $\max \sum_i c_i x_i$，各资源约束写为 $\sum_i a_{ji}x_i\le b_j$。使用单纯形法求解，并逐项复算约束残差。
\section{结果分析}
在基准容量下，方案一的目标值为 120，方案二为 114。相比之下，方案一多获得 6 个收益单位，且三项约束残差均非负，因此采用方案一；其中第二项资源约束取等号，说明它是当前决策的限制因素。
\section{模型检验}
把求解结果回代目标函数和全部约束，复算目标值仍为 120，最大绝对残差为 $2\times10^{-8}$，低于 $10^{-6}$ 的接受阈值。
\section{灵敏度分析}
将第二项容量在基准值上下扰动 10\%，每个扰动点重新优化。容量下降 6\% 以前方案不变，超过该点后第二项决策量开始下降，因此稳定区间为 $[-6\%,10\%]$。
\section{模型评价与改进}
线性结构使约束来源和活跃状态容易核对，但单位收益固定限制了外推范围。后续可用分段线性收益替代常系数，并继续用同一目标口径比较。
\begin{thebibliography}{9}
\bibitem{lp} 运筹学教材编写组. 运筹学[M]. 北京: 高等教育出版社, 2020.
\end{thebibliography}
\appendix
\section{附录}
入口脚本读取资源矩阵，输出决策向量、目标值和逐项约束残差。
\end{document}
